\newglossaryentry{veracity}{
    name={Veracity},
    text={veracity},
    description={Habitual conformity to facts}
}

\newglossaryentry{veracity bond}{
    name={Veracity bond},
    text={veracity bond},
    description={Financial collateral staked by creators to signal confidence in the veracity of their content}
}

\newglossaryentry{counter-veracity bond}{
    name={Counter-veracity bond},
    text={counter-veracity bond},
    description={Financial collateral staked by a challenger when disputing the veracity of content from a creator}
}

\newglossaryentry{misinformation}{
    name={Misinformation},
    text={misinformation},
    description={False information that is spread regardless of intent to deceive as opposed to the intentional deception of disinformation}
}

\newglossaryentry{synthetic content}{
    name={Synthetic content},
    text={synthetic content},
    description={Digital content wholly or substantially generated by non-human systems}
}

\newglossaryentry{prosocial}{
    name={Prosocial},
    text={prosocial},
    description={Someone who is public-spirited and not greedy}
}

\newglossaryentry{creators}{
    name={Creator},
    text={creators},
    description={A participant creating original content}
}

\newglossaryentry{challengers}{
    name={Challenger},
    text={challengers},
    description={Someone questioning the veracity or accuracy of content posted by creators}
}

\newglossaryentry{jurors}{
    name={Juror},
    text={jurors},
    description={An impartial evaluator responsible for adjudicating challenges raised by challengers.}
}

\newglossaryentry{contest}{
    name={Contest},
    text={contest},
    description={A structured dispute in which a creator’s content, backed by a veracity bond, is challenged by one or more challengers staking a counter-veracity bond, with the final outcome determined by a jury’s evaluation of the evidence}
}

\newglossaryentry{challenge period}{
    name={Challenge period},
    text={challenge period},
    description={The allotted period of time for which challenges may dispute the content of a creator}
}

\newglossaryentry{deliberation period}{
    name={Deliberation period},
    text={deliberation period},
    description={The allotted period of time for which jurors may reach a verdict on a particular contest}
}

\newglossaryentry{platform}{
    name={Platform},
    text={platform},
    description={An overarching system or framework that facilitates the interactions among creators, challengers, and jurors. Not to be mistaken for a physical platform}
}
